\documentclass[11pt]{article}
\usepackage[margin=1in]{geometry}
\usepackage{amsmath,amssymb,amsthm}
\usepackage{booktabs,enumitem,array}
\usepackage[hidelinks]{hyperref}
\usepackage[expansion=false]{microtype}
\setlist{nosep,leftmargin=1.5em}

\newtheorem{theorem}{Theorem}
\theoremstyle{remark}
\newtheorem*{remark}{Remark}

\newcommand{\status}[1]{\hfill{\normalfont\small\textsc{[#1]}}}
\newcommand{\one}{\mathbf{1}}
\newcommand{\Wt}{\widetilde{W}}

\title{\textbf{v0.3d patch}\\[2pt]
\large The $n=4$ global ambiguity is now certified; $n\ge5$ is now a theorem}
\author{18 September 2026}
\date{}

\begin{document}
\maketitle
\vspace{-2.5em}

\paragraph{Status.} The simplex bug is real and is fixed; rerunning makes the ambiguity \emph{more} common, not less. The $n=4$ statement has been upgraded from numerical evidence to a certified theorem this session, using a rational witness plus an interval (Krawczyk) certification of the second root, so it no longer waits on a rational pair. Your $n\ge5$ argument is adopted as written, with the aggregate caveat. The DSA feasible-region rerun is clean.

%======================================================================
\section*{1. Domain fixes and rerun}

Weights are now drawn on the simplex (Dirichlet for sampling, softmax with the fourth logit fixed for root search), and roots are accepted only if $w_4=1-w_1-w_2-w_3>0$ together with $\eta_v<p_{2v}<1/2<p_{1v}<1-\eta_v$. For the DSA check, $c_+=\lambda_0\rho$ and $c_-=\lambda_0(1-\rho)$ with $\lambda_0,\rho\in(0,1)$.

\begin{center}\small
\begin{tabular}{@{}lll@{}}
\toprule
Check & Before (box sampling) & After (feasible region)\\
\midrule
Corrupted model, $n=4$ & 7/20 laws with $>1$ root & \textbf{10/20}; root counts up to 3\\
Corrupted model, $n=5$ & 0/10 & 0/12\\
DSA off-atom, $n=4$ & 0/20 & 0/20\\
DSA off-atom, $n=5$ & 0/20 & 0/20\\
Jacobian ranks, $n=4..7$ & full & full (unchanged)\\
\bottomrule
\end{tabular}
\end{center}

So the corrected figure is that roughly half of random in-simplex laws at $n=4$ have several feasible preimages, and the v0.1b DSA claim survives on its correct feasible region.

%======================================================================
\section*{2. $n=4$: global non-identifiability, certified}

\begin{theorem}[$n=4$ is globally non-identifiable on a nonempty open set]\status{proved, computer-assisted}
Let $F$ map $\theta=(w_1,w_2,w_3,p_1,p_2,\eta)$ to the first 15 cell probabilities of the corrupted model at $n=4$. Then there is a nonempty open set $V$ of laws with $|F^{-1}(q)\cap\Theta|\ge2$ for every $q\in V$.
\end{theorem}

\noindent\emph{Witness and certificate.}
\begin{enumerate}
\item $\theta_1$ is rational and strictly feasible:
\[
w=\Bigl(\tfrac1{20},\tfrac3{20},\tfrac3{10}\Bigr),\quad
p_1=\Bigl(\tfrac34,\tfrac{17}{20},\tfrac45,\tfrac{17}{20}\Bigr),\quad
p_2=\Bigl(\tfrac3{10},\tfrac3{20},\tfrac3{20},\tfrac25\Bigr),\quad
\eta=\Bigl(\tfrac3{20},\tfrac1{10},\tfrac1{20},\tfrac1{20}\Bigr),
\]
so $w_4=\tfrac12$ and $\eta_v<p_{2v}<\tfrac12<p_{1v}<1-\eta_v$ for all $v$. The target $p=F(\theta_1)$ is therefore exactly rational, and exact rational elimination gives
\[
\det J(\theta_1)=-\frac{49846858308155727}{655360000000000000000000000000000000}\approx-7.606\times10^{-20}\neq0 .
\]
\item A second root $\theta_2$ of $F(\theta)=p$ was refined to 50 digits ($\max|F(\theta_2)-p|\approx1.7\times10^{-52}$), with $\|\theta_2-\theta_1\|_\infty\approx0.08146$; its weights are $(0.10248,0.17816,0.24636)$ with $w_4\approx0.47300$.
\item \emph{Krawczyk certification.} With $y=\theta_2$, $Y\approx J(y)^{-1}$ and the box $X$ of radius $10^{-6}$ around $y$, interval arithmetic gives $\|I-YJ(X)\|_\infty=5.17\times10^{-2}<1$ and $K(X)\subset\operatorname{int}X$ (the deviation $\max|K-y|$ is $5.17\times10^{-8}$, two orders inside the box). Hence $F(\theta)=p$ has a unique root in $X$, and $J$ is nonsingular on all of $X$. Every point of $X$ satisfies the domain inequalities and $w_4>0$, checked on the box endpoints.
\item $\|\theta_1-\theta_2\|_\infty\approx0.0815\gg10^{-6}$, so the two roots are distinct and both regular. The inverse function theorem gives disjoint neighborhoods on which $F$ is a diffeomorphism onto open images, whose intersection is the required $V$.
\end{enumerate}

This is exactly the route you described, with the second step done by interval certification rather than rational reconstruction --- the second root appears to be algebraic irrational, and chasing rational coordinates was not needed.

%======================================================================
\section*{3. $n\ge5$: generic global identifiability of the aggregate model}

\begin{theorem}[adopted from your argument]\status{proved modulo the standard genericity statements}
For $n\ge5$ the aggregate relation-corrupted model is generically globally identifiable: the four product components and hence $(c_+,c_-,S_+,S_-)$, $\{\tilde\alpha^{\pm}_v\}$ and $\{\eta_v\}$ are determined by the law of $\Wt$, up to a measure-zero exceptional set. The split of $S_+$ into $\lambda_{E+}$ and $\lambda_{B+}$ is \emph{not} determined, at any $n$.
\end{theorem}

The argument, recorded so it can be written up directly: group the views as $G_1=(W_1,W_2)$, $G_2=(W_3,W_4)$, $G_3=(W_5,\dots,W_n)$, with alphabet sizes $4,4,2^{n-4}$. For $r=4$ classes the generic Kruskal ranks are $4,4,\min(4,2^{n-4})$, so at $n=5$ the sum is $4+4+2=10=2r+2$, exactly Kruskal's threshold, and larger $n$ is more permissive; this is the Allman--Matias--Rhodes route for mixtures of product distributions. The tie-constrained submodel does not sit inside the exceptional set: at the rational witness the two grouped determinants are $\det M_{12}=-63/8000$ and $\det M_{34}=-2783/116640$, both nonzero, and a fifth view with four pairwise-distinct component probabilities, e.g.\ $(p_{1,5},p_{2,5},1-\eta_5,\eta_5)=(3/4,1/4,9/10,1/10)$, has Kruskal rank 2. Labels are then fixed: generically only one pair of components satisfies $q_{a,v}+q_{b,v}=1$ for all $v$ (the atom pair), orientation $q_{A^+,v}>1/2>q_{A^-,v}$ separates $A^+$ from $A^-$, and $p_1>1/2>p_2$ separates $D^+$ from $D^-$. Finally $\alpha^+_v=\tfrac12+(p_{1v}-\tfrac12)/(1-2\eta_v)$ and $\alpha^-_v=\tfrac12+(\tfrac12-p_{2v})/(1-2\eta_v)$.

\begin{remark}
Stating the aggregate caveat prominently is important: what more views buy is the latent verifier-and-check structure, never the easy-correct versus blind-wrong equivalence, which is Corollary~2 and needs gold labels. The two results should sit next to each other for exactly that contrast.
\end{remark}

%======================================================================
\section*{4. The statement to carry into the paper}

\[
\boxed{\begin{array}{ll}
n\le3 &: \text{not locally identifiable, by dimension;}\\[1mm]
n=4 &: \text{generically locally identifiable, globally ambiguous on an open set (certified);}\\[1mm]
n\ge5 &: \text{generically globally identifiable for the aggregate corrupted model;}\\[1mm]
\text{all }n &: \lambda_{E+}\ \text{vs.}\ \lambda_{B+}\ \text{is not identified without gold labels.}
\end{array}}
\]

%======================================================================
\section*{5. Remaining item, and one terminology fix}

The open piece is now the generic fiber degree at $n=4$. $F$ is dominant and generically finite, so away from the discriminant the complex fiber has a fixed degree $d$; we have a certified $d\ge2$ and have observed up to three feasible real roots (you saw four), so $d$ is likely larger still. Monodromy or homotopy continuation would settle it, and the resulting statement --- $n=4$ is generically $d$-to-one, with the number of \emph{feasible real} roots varying by region --- is a cleaner explanation of the scan than any list of counts. That needs a solver we do not have here.

Terminology: the elimination in \texttt{global\_id.py} is exact rational Gaussian elimination with \texttt{Fraction}, not fraction-free Bareiss elimination. Corrected in the script comments and above.

\paragraph{Scripts.}\sloppy \texttt{scan2.py} (simplex-constrained scan), \texttt{dsa\_global.py} (feasible-region DSA check), \texttt{witness.py} (search for a rational-target double root), \texttt{certify.py} (50-digit refinement, Krawczyk certification, exact determinant, box feasibility), \texttt{corrupted.py} (Dirichlet weights, ranks and conditioning).

\end{document}